\section{A Choir of Witnesses: The Living Ark in the Church's Memory}

The Fathers did not begin with an ingenious resemblance and then promote it into
doctrine. They began with the identity of Jesus Christ. The child whom Mary
bears is the eternal Word, the Lawgiver who once spoke above the mercy seat, the
true manna, and Emmanuel. Once that Christological center is held fast, the Ark
becomes more than an ornamental title. The old chest carried signs of God's
covenant presence; Mary carried, according to his humanity, the divine person to
whom every sign pointed. The Church's teachers and worship therefore sound a
common refrain across Greek, Syriac, Latin, Byzantine, and Roman voices: the
type was holy, but the living fulfillment is holier because the Lord himself
dwells within her.

This chorus is not flattened into a claim that every Father used the same title
or drew every later conclusion. A genuine homily, a work preserved in
translation, an anonymous hymn, a medieval sermon, and an official modern
liturgical text remain different kinds of witnesses. Heard in their proper
parts, however, they disclose a tradition both ancient and cumulative.

\subsection{The living Word within the living Ark}

The earliest secure direct witness in this study comes from St.~Ephrem the
Syrian. In the sixteenth \emph{Hymn on the Nativity}, Ephrem pictures Joseph
rising to minister before his Lord dwelling in Mary, as a priest ministers
before the Ark because of divine holiness. The next stanza contrasts Moses,
who carried stone tablets, with Joseph attending the pure tablet in whom the
Creator's Son dwells (\emph{Hymns on the Nativity} 16.16--17). The image is
already domestic and Christological: Mary is the Ark because the Son is within
her, and Joseph's service is active attendance upon the Lord and his mother.
Elsewhere Ephrem calls Mary ``that Ark'' and contemplates the hidden power
within her overthrowing Satan as the captured Ark overthrew Dagon
(\emph{Hymns on the Nativity} 4.112--13). The Ark is victorious not by Mary's
autonomous power but because the Lord has entered his living sanctuary.

St.~Epiphanius of Salamis soon makes the biblical correspondence still more
explicit. His Greek \emph{On Weights and Measures}, written in 392, survives
completely in Syriac at the relevant passage. In section 35 he calls Mary the
holy or living Ark prefigured by the wilderness Ark. The former contained the
written word on stone; Mary bore the living Word. Epiphanius then reads the
Davidic correspondence: when the living Ark entered Elizabeth's house, John
danced within his mother's womb as David had danced before the Ark. This is not
a modern concordance imposed upon two stories. It is a fourth-century Father
reading the Visitation within Israel's sacred history.

Across the same late-fourth- or early-fifth-century horizon, the Latin West
answers with St.~Maximus of Turin. In a sermon now classed as most probably
genuine, he begins from David's dance and asks what the Ark can mean if not holy
Mary. The old Ark contained tablets, Law, and the divine voice; Mary bears the
covenant's heir, Gospel, and true Word. Gold within and without becomes the
splendor of virginity in the inward and outward life (\emph{Sermon} 42.5).
Older printings placed this sermon under St.~Ambrose. Naming Maximus restores
rather than diminishes the early Latin witness.

Early in the fifth century St.~Proclus of Constantinople gives the image a
concentrated Christological form. In his authentic fifth homily he venerates
Mary as mother, servant, cloud, bridal chamber, and ``Ark of the Lord.'' He
immediately controls the comparison: not an ark carrying tablets, but the woman
who carried the Lawgiver. Cloud, chamber, and Ark do not compete. Together they
confess that the transcendent Son truly entered a human history and a human
womb without ceasing to be Lord.

St.~Hesychius of Jerusalem expands the same confession soon after the Council
of Ephesus. Preaching on the Mother of God, he applies Psalm 132:8---``Arise,
O Lord, into your rest, you and the ark of your holiness''---to the Virgin
\emph{Theotokos}. He also sets Mary's Ark beside Noah's: the ancient vessel
preserved mortal life from the flood, while Mary bore incorruptible Life and
the Creator of Noah. His trinitarian imagery is especially important for this
study. The~Spirit comes upon her, the Father's power overshadows her, and the
Son dwells within her. The Ark title consequently belongs inside the mystery of
the missions of the Son and Spirit; it never makes Mary a source of divinity.

Later in the same Jerusalem preaching world, a homily transmitted under the
name of the monk Chrysippus calls the Virgin a royal and most precious Ark. It
distinguishes her from Noah's Ark and from the chest that carried stone tablets:
the Creator is both her maker and her indwelling Lord. Jugie defended the
attribution by comparison with Chrysippus's secure preaching, but the principal
surviving manuscript is eleventh-century and this study did not independently
establish a modern critical consensus. The text therefore remains a strong,
qualified reception witness traditionally assigned to a fifth-century setting,
rather than an unqualified authenticity anchor.

\begin{fourstable}{Witness}{Ark confession}{Christological center}{Proper limit}
Ephrem, fourth century & Joseph ministers before the Lord dwelling in Mary as a priest before the Ark. & Stone tablets yield to the pure tablet in whom the Creator's Son dwells. & The prose is an independent paraphrase anchored to Beck's identified critical locus and checked translations; no direct English quotation is claimed.\\
Epiphanius, 392 & Mary is the holy or living Ark; John dances before her. & The written word yields to the living Word whom she bears. & The passage is quoted through the complete Syriac version of a Greek work.\\
Maximus of Turin, late fourth or early fifth century & David dances before the Ark that signifies holy Mary. & Tablets, Law, and voice yield to heir, Gospel, and true Word. & Cite the modern Maximus corpus, not the obsolete Ambrosian attribution.\\
Proclus, early fifth century & Mary is mother, cloud, chamber, and Ark of the Lord. & She carries the Lawgiver, not merely tablets. & The title does not confuse servant and Lord.\\
Hesychius, fifth century & The Virgin is the Ark of Psalm 132 and bears incorruptible Life. & Father, Son, and Spirit order the mystery of divine indwelling. & Psalmic application is typological, not the psalm's sole literal sense.\\
Chrysippus transmission, traditionally fifth century & Mary is the royal and precious Ark. & Her maker is the Creator who chooses to dwell within her. & Jugie defended the transmitted attribution by stylistic comparison; the principal surviving manuscript is eleventh-century.\\
\end{fourstable}

\subsection{Gold, incorruptible wood, and the poetry of holiness}

The material features of the Ark invited contemplative reading. Gold within
and without suggested an integrity embracing the hidden and visible life.
The Septuagint's expression for the Ark's acacia boards, ``incorruptible
wood,'' made the vessel a ready image of virginal integrity and a body wholly
consecrated to God. Such readings are neither claims about the botanical
indestructibility of acacia nor early verbatim definitions of the Immaculate
Conception. They are the Church's language of fittingness: the all-holy Son
prepared and sanctified the mother from whom he took flesh.

An attribution must sometimes be surrendered so that a text can speak truthfully.
A homily formerly printed under Gregory Thaumaturgus magnificently calls the
Virgin an Ark gilded within and without, but it is not secure third-century
Gregory. Similar rich texts circulated under Athanasius, Chrysostom, Proclus,
Ildefonsus, and Bernard without belonging securely to those authors. Their
images remain evidence of Christian reception when labeled as pseudonymous;
they may not be used to manufacture an earlier date or a famous endorsement.
The faith loses nothing when the choir's anonymous singers are allowed their
own voices.

St.~Bonaventure supplies a secure but more focused medieval witness. In an
authentic sermon for the Purification he identifies the Ark overshadowed by the
cherubim with the glorious Virgin, filled with divine light and wholly intent
upon divine things. In his explicit mapping, the gold signifies contemplative
purity (\emph{On the Purification of the Blessed Virgin}, Sermo I,
\emph{collatio}). A much richer thirteenth-century sermon printed in the older
Quaracchi Bonaventure corpus organizes the Ark by material, contents, efficacy,
and honor and asks whether virginity is not a kind of incorruptibility. Modern
critical work did not retain that sermon among Bonaventure's secure works. Its
theology remains valuable medieval reception, but the richer synthesis is not
placed on the saint's authority.

Other secure medieval voices broaden the chorus. Peter of Celle places Mary's
betrothal to Joseph, the Spirit's overshadowing, divine birth, temple, and Ark
within one chapter of his \emph{Book of Breads}; the mother is the Ark of the
Covenant because she contains within herself the sanctification of creature
and Creator. St.~Anthony of Padua
then applies the Ark's arising to Mary's bodily Assumption into the heavenly
bridal chamber. Betrothal, indwelling, and heavenly rest belong to one
Christ-centered arc without collapsing into one proof.

\subsection{The Ark enters the heavenly sanctuary}

The tradition reaches a luminous maturity in St.~John Damascene's authentic
homilies on the Dormition. He calls Mary the spiritual and living Ark of God
the Word and contemplates her entry into the heavenly Holy of Holies. The image
has moved from Sinai, procession, and Temple to eschatological fulfillment:
the body from which the Word took flesh does not cease to belong to his victory.
Here Psalm 132:8, the Ark's ascent, and the Church's Dormition faith converge
as a powerful fittingness for the Assumption.

The Byzantine liturgy makes that reception communal. The Akathist greets Mary
as an Ark gilded by the Spirit. The Annunciation canon addresses her as the
living Ark of God; the Dormition canons gather Ark, lampstand, manna, Aaron's
rod, tablet, tabernacle, and Holy of Holies around the one mystery of the Word
made flesh. Poetic abundance does not blur the center: Christ is the treasure,
Mary the creature who receives and bears him. Nor does a liturgical image by
itself define a dogma. It witnesses the Church praying within the typology.

The Roman Church sings the same faith in another register. The Litany of Loreto
invokes Mary as \emph{Foederis arca}, ``Ark of the Covenant.'' The official
\emph{Collection of Masses of the Blessed Virgin Mary} is still more explicit.
Its Visitation formulary calls Mary the Ark of the New Covenant who brings
salvation and joy into Elizabeth's house; its Mass of Mary, Temple of the Lord,
calls her the Ark containing Jesus Christ, author of the new law. Ancient
exegesis has become the Church's present worship without ceasing to be
Christological.

The modern Magisterium deliberately receives this inheritance. St.~Paul VI
calls Mary the true Ark and true Temple because at the Annunciation she received
the one Mediator in her body; he links the Spirit's overshadowing, consecrated
virginity, and the Fathers' Ark titles (\emph{Marialis cultus} 6, 26). The
\emph{Catechism} names her Daughter Zion in person, Ark of the Covenant, and
dwelling of the Lord's glory (CCC 2676). Benedict XVI then reads the Visitation
as the new Ark's procession: Elizabeth recognizes the Lord's presence and the
child's joy recalls David. These acts do not invent a new typology. They show
the Church receiving a patristic and liturgical reading as living catechesis.

\subsection{Joseph before the mystery entrusted to him}

The Ark tradition must illuminate Joseph and Mary's marriage positively. The
Gospel does not cast Joseph as another Uzzah standing before a dangerous woman.
God commands him to take Mary as his wife. He obeys, receives her into his home,
names the child, protects mother and Son, and establishes the Davidic household
in which the incarnate Word is welcomed. His continence is not recoil from the
body or contempt for marriage; it is the freely lived form of this singular,
real marital vocation.

The Catholic tradition therefore places reverence inside covenantal love.
Origen calls Joseph the steward or minister of the Lord's birth
(\emph{Homilies on Luke} 13.7). St.~Ambrose insists that the name ``spouse''
does not take virginity away but bears witness to a marriage
(\emph{Exposition of Luke} II.5). St.~Augustine supplies the fullest patristic
account: freely chosen perpetual continence makes the bond firmer rather than
breaking it; Mary and Joseph are spouses and parents in mind and purpose though
not through Joseph's begetting (\emph{On Marriage and Concupiscence}
I.11.12--13). Chastity is the form of their fidelity, not the negation of their
marital gift.

St.~Ephrem had already given that vocation a liturgical and domestic shape in
Joseph's service before the Ark. St.~Bernard of Clairvaux gives the same service
an affectionate inventory. The mystery is entrusted to Joseph, who is permitted
not only to see and hear Christ but to
carry, lead, embrace, kiss, nourish, and guard him (\emph{In Praise of the
Virgin Mother} II.16). St.~Thomas Aquinas gathers the doctrine: this marriage
was simply true, founded upon an inseparable union of souls and mutual consent,
and brought to its singular perfection in the upbringing of Christ (ST III,
q.29, a.2). Husband and wife receive together, in distinct personal ways, the
Son who chose to be entrusted to them.

The Roman liturgy presents their bond as a love at once marital and virginal,
and St.~John Paul II describes their exclusive gift of self as conjugal love born
anew in the Holy Spirit (\emph{Redemptoris custos} 18--20). Joseph's reverence
is therefore active and tender: he receives, names, carries, works, journeys,
protects, nourishes, and gives himself. He does not seize a sacred object. He
loves his wife and serves the embodied Lord whom she bears.

Ephrem and Matthew supply the governing positive axis. Service once rendered
before the Ark reaches its personal and covenantal fulfillment: Joseph obeys
the command to take Mary as his wife, receives the Lord dwelling in her, and
gives himself to the guardianship of their household. The fulfillment belongs
to appointed Ark-service becoming spousal and fatherly communion; it does not
belong to Uzzah's death. Uzzah remains only a bounded contrast from a
disordered procession. He does not prefigure Joseph and does not explain the
marriage. Marital union is not declared unclean, Mary is not imagined as
physically dangerous, and Joseph's chastity is not reduced to fear.

\interpretivekey{ordered service becomes covenantal guardianship}{Ephrem's priest-before-the-Ark image and Matthew's divine command provide the controlling pattern: Joseph takes Mary as wife, receives the incarnate Lord, and gives himself in faithful, virginal guardianship. Uzzah is only a contrast concerning disordered handling, never a type fulfilled by Joseph or a reason to regard marriage as defiling.}

\subsection{One faith, many registers}

The witnesses do not all prove the same proposition. Ephrem, Epiphanius,
Maximus, and the Jerusalem homilists establish ancient explicit Ark reception.
Proclus protects its Christological center. Byzantine hymnography shows the
image becoming the Church's prayer. John Damascene unfolds its eschatological horizon.
Bonaventure contemplates the Ark's golden purity, while broader medieval
reception unfolds its material and contents. Roman liturgy receives it
as a public title and makes its Visitation setting explicit. Together they do
not replace Scripture, the definitions of the Church, or the proper arguments
for the four Marian dogmas. They show that the New-Ark reading is no isolated
modern conceit but a durable Catholic way of confessing the Incarnation.

\sourceboundary{Authenticated patristic and saintly witnesses, honestly labeled traditional texts, and received Byzantine and Roman worship}{An ancient and enduring Christological reception of Mary as the living Ark; a secure David--John reading by 392; the Lawgiver/tablets contrast; holiness, virginity, and heavenly-sanctuary fittingness; positive virginal guardianship in Joseph and Mary's true marriage}{A claim that every Father taught the typology, that a pseudonymous work belongs to its traditional namesake, that Ark imagery independently defines all four dogmas, or that Joseph is the New Testament Uzzah}
